\section{Advanced Robotics \& Automation}\label{sec:automation}

Automation has long underpinned industrial production, but its nature has shifted dramatically with the transition from rigid mechanization to intelligent, connected systems.
Early automation relied on fixed robotic arms, conveyor systems, and programmable logic controllers (PLCs), which were efficient for mass production but inflexible in the face of product variability~\cite{brettel-how-virtualization-decentralization-and-networking, monsori-cpps-roots-expectations-rnd-challenges}.
In \ac{hmlv} manufacturing, this rigidity proved especially limiting, as frequent changeovers, short product life cycles, and small batch sizes demand adaptability.

\subsection{From Fixed Automation to Flexible Systems}

Industry 4.0 brought a new paradigm by embedding robotics into broader cyber-physical production systems~\cite{monsori-cps-in-manufacturing}.
Robots are no longer isolated automation units, but interconnected agents within systems of sensing, data exchange, and decision-making.
This has enabled a shift toward more flexible, responsive forms of automation.

One important development has been the rise of collaborative robots (cobots), which integrate advanced sensing, force control, and safety mechanisms to share workspaces with human operators~\cite{villani-survey-on-human-robot-collaboration-in-industrial-settings}.
Their flexibility makes them particularly relevant to \ac{hmlv} production, where some assembly or finishing tasks require human dexterity while still benefiting from robotic consistency.
IN parallel, reconfigurable robotic systems have merged, emphasizing modular hardware and software architectures that can be adapted to different product types or processes with minimal downtime~\cite{Bi-reconfigurable-manufacturing-systems-state-of-the-art,bortolini-reconfigurable-manufacturing-systems-literature-review}.
While the concept of cobots leans more towards the human-centric ideals of "Industry 5.0", the improvement in quality, efficiency and partial automation also makes it relevant to the Industry 4.0 industrial paradigm.

\subsection{Intelligence and Adaptivity in Automation}

Advances in artificial intelligence and machine learning extend the scope of robotics beyond repetitive automation.
Machine vision, reinforcement learning, and adaptive motion planning now allows robots to handle semi-structured environments and uncertain inputs~\cite{kragic-interactive-collaborative-robots-challenges-and}.
These capabilities are especially valuable in \ac{hmlv} contexts, where production tasks vary and cannot be exhaustively pre-programed.
Such developments are closely tied to big data analytics and \ac{iot} infrastructures~\cite{LEE-2015-a-cps-for-indsutry40-based-manufacturing-systems,xu-industry-40-state-of-the-art-and-future-trends}, since adaptive automation depends on streams of sensor data and predictive models for decision support.

\subsection{Networked and Cloud Robotics}

The rise of cloud robotics and edge-assisted automation illustrates another step in this evolution.
By linking robotic systems through networked infrastructures, capabilities such as shared learning, fleet coordination, or remote monitoring become feasible~\cite{kehow-a-survey-of-research-on-cloud-robotics-and-automation}.
These approaches align with the broader Industry 4.0 emphasis on connectivity and data-driven optimization, but also raise questions of latency, security, and interoperability.

\subsection{Challenges and Future Directions}

Despite rapid progress, challenges remain. Advanced robotics often entails high upfront costs, and the return on investment can be difficult to justify in small-back production unless flexibility and reusability are maximized~\cite{bortolini-reconfigurable-manufacturing-systems-literature-review}.
Human-robot collaboration also introduces ergonomic, trust, and safety considerations that are not yet fully resolved~\cite{villani-survey-on-human-robot-collaboration-in-industrial-settings}.
Integration into wider CPS and \ac{iot} infrastructures is further slowed by the lack of universal standards, a theme which is common across Industry 4.0 technologies~\cite{monsori-cpps-roots-expectations-rnd-challenges,xu-industry-40-state-of-the-art-and-future-trends}.

Overall, robotics and automation are evolving from rigid, standalone systems into flexible, intelligent, and networked production agents. Their ongoing development is central to achieving the adaptability and responsiveness for \ac{hmlv} manufacturing, even as integration, cost, and safety challenges remain active areas of research and industrial practice.
